% mage-auto-heading: name-the-capability
% mage-auto-heading: enumerate-the-failure-classes
% mage-auto-heading: assign-one-guarantee-to-each-failure
% mage-auto-heading: find-the-dependency-among-guarantees
% mage-auto-heading: separate-the-load-bearing-path-from-the-enhancements
% mage-auto-heading: match-the-substrate
% mage-auto-heading: stop
% mage-auto-heading: the-method-in-one-picture
% Authoritative LaTeX manuscript; Markdown is a generated compatibility projection.

The seven preceding stacks are examples. The reusable method is how to compose one. Start from the capability, not the mechanisms.

\section{1. Name the Capability}\label{heading:name-the-capability}

Prefer

% <!-- inline-quote -->

\begin{quote}
Every consequential mutation is reconstructable.
\end{quote}

over

% <!-- inline-quote -->

\begin{quote}
We need provenance logging.
\end{quote}

The first states the engineering result. The second has already chosen a mechanism, before you know whether it is the right one or the only one.

\section{2. Enumerate the Failure Classes}\label{heading:enumerate-the-failure-classes}

Ask what would make the capability claim false. For reconstructable mutation:

\begin{itemize}
\tightlist
\item
  the actor is unknown;
\item
  a record is missing;
\item
  a bypass path exists;
\item
  history cannot be joined;
\item
  the transformation silently loses semantics.
\end{itemize}

Those failures determine what the composition must cover.

\section{3. Assign One Guarantee to Each Failure}\label{heading:assign-one-guarantee-to-each-failure}

Map each way the claim can break to the guarantee that closes it.

{\footnotesize
{\def\LTcaptype{none} % do not increment counter
\begin{longtable}[]{@{}ll@{}}
\toprule\noalign{}
Failure & Guarantee \\
\midrule\noalign{}
\endhead
\bottomrule\noalign{}
\endlastfoot
actor unknown & attribution \\
record missing & emission \\
bypass exists & coverage \\
history fragmented & stable identity / join \\
semantic loss & fidelity check \\
\end{longtable}
}
}

Do not add a mechanism because the case you are copying from used it. Add it because a failure class needs its guarantee.

\section{4. Find the Dependency Among Guarantees}\label{heading:find-the-dependency-among-guarantees}

Some guarantees only become meaningful once another exists.

\begin{verbatim}
MARK  →  EMIT  →  COVER  →  READ
\end{verbatim}

A completeness check is meaningful only after the population to be accounted for has been defined. A list of mechanisms becomes a stack when their guarantees depend on one another.

\section{5. Separate the Load-Bearing Path from the Enhancements}\label{heading:separate-the-load-bearing-path-from-the-enhancements}

For each mechanism, ask:

% <!-- inline-quote -->

\begin{quote}
If I remove this, is the capability claim still valid?
\end{quote}

If yes, the mechanism may still be valuable, but it is not load-bearing. The diagrams represent this distinction with solid paths and dashed attachments.

\section{6. Match the Substrate}\label{heading:match-the-substrate}

Do not reproduce a reference stack mechanically. A database, an embedded controller, a compiler, a SaaS application, and a research codebase each expose different observable facts and different deterministic seams. Preserve the guarantees; let the substrate decide the mechanisms.

\section{7. Stop}\label{heading:stop}

The objective is not maximum governance. Stop when the smallest required stack makes the capability claim true for the failure classes that matter.

\section{The Method in One Picture}\label{heading:the-method-in-one-picture}

\begin{figure}[htbp]
\centering
\includegraphics[width=\linewidth]{book/latex/assets/appendix-a-composing-a-stack-b82cb0c3c4.pdf}
\caption{Composing an engineering stack. Begin with a capability and its failure classes. Identify the guarantees required to close those failures, select mechanisms that provide them, determine their dependencies, and retain the smallest required composition that makes the capability claim valid.}
\label{appendix-a-composing-a-stack}
\end{figure}
